\documentclass[twocolumn]{iopjournal}

\usepackage{amsmath,amssymb,amsfonts}

\fancyhfoffset[L]{0mm}
\setlength{\marginparwidth}{0pt}
\setlength{\marginparsep}{0pt}

% ==========================================
% ZENODO PREPRINT HEADERS (ACTIVE)
% ==========================================
\renewcommand{\articletype}[1]{%
  {\vspace*{-8mm}\noindent \Large \sf Preprint}\par
  \vspace*{6mm}\noindent{\scriptsize \sf{\bfseries \MakeUppercase{#1}}}}

\fancyhead[L]{}
\fancyhead[R]{}

\hypersetup{pdftitle={Supplementary Material: The Tetrad Scheme and the Two Bounded Regions}, pdfauthor={Stephen Nelson}, pdfsubject={Gravity Field Dynamics (GFD)}, pdfkeywords={tetrad, scheme, algebraic geometry, Lorentz group, spatial parity, root-blindness, stellated octahedron, proof}}

\begin{document}

\articletype{Supplementary Material}

\title{Supplementary Material: The Tetrad Scheme\\and the Two Bounded Regions}

\author{Stephen Nelson}

\affil{Field Dynamics Research}

\begin{abstract}
The main letter derives the stellated octahedron from the half-space inequalities of four derived tetrahedral directions (Route~1). This note provides a second, independent route to the same result (Route~2). The defining equation of the tetrad, $g_{\mu\nu} = \eta_{ab}\, e^a{}_\mu\, e^b{}_\nu$, is quadratic in the tetrad, and it forces $(\det e)^2 = \det g / \det \eta$, an equation with two roots. Its solution space (scheme) is two schemes, one over each root, related by spatial parity. Each carries a regular tetrahedron with a $3+1$ face decomposition. Their compound is the stellated octahedron. The metric determines only $(\det e)^2$, so it cannot distinguish them. This argument uses no half-space bounding, no recession cones, and no sign-combination enumeration.
\end{abstract}

\keywords{tetrad, scheme, algebraic geometry, Lorentz group, spatial parity, root-blindness, stellated octahedron, proof}

% ----------------------------------------------------------------------

\section*{The Tetrad Equation Is Quadratic}

The tetrad $e^a{}_\mu$ is defined by the relation
\begin{equation}
g_{\mu\nu} = \eta_{ab}\, e^a{}_\mu\, e^b{}_\nu, \qquad\text{or}\qquad g = e^T \eta\, e,
\label{eq:tetrad}
\end{equation}
which is quadratic in $e$. If $e^a{}_\mu$ is a solution, then so is $\Lambda^a{}_b\, e^b{}_\mu$ for any $\Lambda \in O(1,3)$, because $O(1,3)$ preserves $\eta_{ab}$ by definition. Conversely, every solution is of this form. Taking determinants in Eq.~\eqref{eq:tetrad} gives $(\det e)^2 \det \eta = \det g \neq 0$, so every solution is invertible; for two solutions $e$ and $e'$, the matrix $\Lambda = e' e^{-1}$ satisfies $\Lambda^T \eta \Lambda = e^{-T} g\, e^{-1} = \eta$. The full solution space at any point is therefore the orbit of a single tetrad under the Lorentz group $O(1,3)$, and distinct $\Lambda$ give distinct solutions.

This solution space is what algebraic geometry calls the \emph{scheme} of the equation~\cite{Grothendieck1960,Hartshorne1977}: the set of all solutions, considered as a geometric object in its own right, without selecting a preferred root. Precisely, the sixteen entries of $e$ are the unknowns, Eq.~\eqref{eq:tetrad} is a set of quadratic equations in them, and
\[
X_g = \operatorname{Spec} \mathbb{R}[e^a{}_\mu] \,/\, (e^T \eta\, e - g)
\]
is the affine scheme they define. Its real points are the tetrads, and by the same algebra, which holds over any ring, $X_g$ is isomorphic to the orthogonal group scheme $O(\eta)$, through $\Lambda \mapsto \Lambda e_0$ for any one tetrad $e_0$.

% ----------------------------------------------------------------------

\section*{The Scheme Is Two Schemes}

Taking determinants in Eq.~\eqref{eq:tetrad} once more,
\begin{equation}
(\det e)^2 = \frac{\det g}{\det \eta} \equiv \kappa > 0,
\label{eq:roots}
\end{equation}
where $\kappa > 0$ because $g$ and $\eta$ have the same signature. This is $y^2 = \kappa$, with roots $\pm\sqrt{\kappa}$, as $x^2 = a$ has $\pm\sqrt{a}$. Every tetrad has one root, never both, so the scheme is two schemes, $X_g = X_+ \sqcup X_-$, one over each root, each open and closed in $X_g$. Neither splits further: under $X_g \cong O(\eta)$ they are $SO(\eta)$ and its coset, and $SO(\eta)$ is connected even over $\mathbb{C}$, where it is $SO(4, \mathbb{C})$~\cite{Borel1991}.

The Lorentz group $O(1,3)$ has four connected components~\cite{Wald1984}, distinguished by two discrete operations:
\begin{itemize}
\item spatial parity $P$: $e^0 \to e^0$, $e^i \to -e^i$ ($i = 1,2,3$),
\item time reversal $T$: $e^0 \to -e^0$, $e^i \to e^i$.
\end{itemize}
The four components are $SO^+(1,3)$, $P \cdot SO^+(1,3)$, $T \cdot SO^+(1,3)$, and $PT \cdot SO^+(1,3)$. A change of frame $e \to \Lambda e$ multiplies $\det e$ by $\det \Lambda$, so it exchanges the two schemes if and only if $\det \Lambda = -1$. $P$ and $T$ each have determinant $-1$ and exchange them; $PT = -I$ has determinant $+1$ and does not. The four real components of $X_g$ are the two schemes, each split by the direction of the time leg.

The main letter's construction uses only the spatial legs of the tetrad: the four derived directions $v_1, v_2, v_3, v_4$ are linear combinations of $e^1, e^2, e^3$. Time reversal $T$ acts only on $e^0$ and leaves the spatial legs unchanged. To compare the frames of one observer, hold the time leg fixed and future-pointing. The changes of frame that do so are $\Lambda = \mathrm{diag}(1, Q)$ with $Q \in O(3)$, and $\det \Lambda = \det Q$: for these frames the root of Eq.~\eqref{eq:roots} is the orientation of the spatial legs.

For one observer, the two schemes are thus the two components of $O(3)$:
\begin{itemize}
\item the identity component $SO(3)$, which rotates the four derived directions rigidly,
\item the parity component $P \cdot SO(3)$, containing $P$, which negates all three spatial legs.
\end{itemize}

The continuous part, $SO(3)$, rotates the four derived directions rigidly and produces no new geometry. The discrete part, $P$, sends every $v_i \to -v_i$, whose normalised directions are a second set of four unit vectors with the same pairwise inner product $c = -1/3$. Four unit vectors in $\mathbb{R}^3$ with common pairwise inner product $-1/3$ form the vertices of a regular tetrahedron~\cite{Coxeter1973}; their convex hull is a bounded solid. The identity component gives TetA. The parity component gives TetB. Precisely, $P$ maps $\mathrm{TetA} = \{\max_i v_i \le 1\}$ to $\{\max_i (-v_i) \le 1\}$, the all-flipped sign combination, which is the second bounded region of Route~1.

Three facts make this identification exact. First, $P$ is the canonical member of its component: the only elements of $O(3)$ commuting with every rotation and reflection are $\pm I$, since commuting with all diagonal sign matrices and all permutation matrices forces a multiple of $I$, hence $\pm I$, and of these only $-I$ has determinant $-1$. Second, parity is not a rotation in disguise: the four directions span $\mathbb{R}^3$, so $Q v_i = -v_i$ for all $i$ forces $Q = -I$, and no element of $SO(3)$ carries the frame of TetA to that of TetB leg by leg. Third, this rests on the spatial dimension being odd: $\det(-I) = (-1)^D$ in $D$ spatial dimensions, and in an even number of spatial dimensions $-I$ is a rotation, so parity would produce no new component.

Two cautions keep the statement exact. As unlabelled point sets the tetrahedra are congruent, since a regular tetrahedron is its own mirror image (a $90^\circ$ rotation about one axis maps the vertices of TetA onto those of TetB), so the two schemes differ in the order of the legs and the region of space they occupy, not in shape. And reversing every leg, $e \to -e = PT\,e$, keeps $\det e$ but also reverses the time leg; its spatial legs are those of $Pe$. The identification of root with tetrahedron therefore holds for future-pointing frames, the setting of the main letter.

% ----------------------------------------------------------------------

\section*{Both Schemes Are Necessarily Present}

The scheme of a polynomial equation contains all its roots by definition. The quadratic structure of the tetrad equation and the two schemes it splits into together guarantee that both TetA and TetB exist in the solution space. Both satisfy the same defining equation with the same metric, and the metric cannot tell them apart: by Eq.~\eqref{eq:roots} it determines $(\det e)^2$ and never $\det e$, and every quantity built from $g$, including the Christoffel symbols, the curvature, the Einstein tensor and the Einstein--Hilbert action, takes the same value on both. Retaining one scheme is consistent, but it replaces the tetrad equation with the system $g = e^T \eta\, e$, $\det e = +\sqrt{\kappa}$, whose second equation $g$ does not contain, just as $x = +\sqrt{a}$ is not $x^2 = a$.

The textbook convention of restricting local frames to the proper orthochronous Lorentz group $SO^+(1,3)$~\cite{Wald1984,Hehl1976} corresponds to working within one of the two schemes. This is a valid computational choice, but it is a restriction of the full scheme to one of them. The scheme itself contains both. The choice between the signature conventions $(-,+,+,+)$ and $(+,-,-,-)$ is a different matter: it changes the sign of $\eta$ as a whole, leaves the Lorentz group unchanged, and selects neither.

% ----------------------------------------------------------------------

\section*{Each Scheme Carries a $3+1$ Decomposition}

In the main letter's coordinates, whose axes are the frame's spatial legs, the four face normals of TetA, expressed as sign patterns on the three spatial axes, are:
\begin{equation}
(+,+,+), \quad (+,-,-), \quad (-,+,-), \quad (-,-,+).
\end{equation}
One face normal has uniform sign: $(+,+,+)$, the normal of $v_1$. Three face normals have mixed signs. This is the $3+1$ partition of the four faces that the main letter reads as one temporal and three spatial. The partition is fixed by the frame's spatial legs: under the tetrahedron's own symmetries all four faces are equivalent, so it is a property of the tetrahedron relative to the frame, not of the tetrahedron alone.

Under spatial parity $P$, the face normals of TetB are:
\begin{equation}
(-,-,-), \quad (-,+,+), \quad (+,-,+), \quad (+,+,-).
\end{equation}
Again, one face normal has uniform sign: $(-,-,-)$. Three have mixed signs. TetB carries its own $3+1$ decomposition, relative to the same legs.

The stellated octahedron therefore carries the $3+1$ pattern in both tetrahedra, with the uniform-sign normals $(+,+,+)$ and $(-,-,-)$ exchanged by spatial parity, independently in both bounded regions.

% ----------------------------------------------------------------------

\section*{Summary}

The main letter derives the stellated octahedron by classifying the sixteen half-space sign combinations and proving exactly two yield bounded regions (Route~1). The present note derives the same result from the quadratic structure of the tetrad equation (Route~2):

\vspace{6pt}
\noindent
{\footnotesize\setlength{\tabcolsep}{4pt}%
\begin{tabular*}{\columnwidth}{@{\extracolsep{\fill}}lll@{}}
\hline
 & Route 1 (letter) & Route 2 (this note) \\
\hline
Input & half-space inequalities & tetrad eqn.\ $g = \eta\, e\, e$ \\
Method & recession cone test & scheme of quad.\ eqn. \\
Key fact & $\sum v_i = 0$ (normals & $(\det e)^2 = \kappa$ has \\
 & sum to zero) & two roots \\
Output & 2 of 16 bounded & 2 schemes \\
\hline
\end{tabular*}}
\vspace{6pt}

Both routes arrive at the same stellated octahedron with the same $3+1$ face decomposition in each tetrahedron. Route~2 uses no half-space construction and Route~1 no quadratic structure; both take the four tetrahedral directions from the main letter's signature-to-frame supplement. A rejection of one leaves the other intact. The loop between equation and geometry is closed: the two roots of the tetrad equation are the two bounded regions, TetA and TetB, exchanged by spatial parity. The metric is blind to which root a frame occupies, and keeping one is a convention.

% ----------------------------------------------------------------------

\begin{thebibliography}{99}

\bibitem{Grothendieck1960}
A.~Grothendieck and J.~Dieudonn\'{e},
\textit{\'{E}l\'{e}ments de g\'{e}om\'{e}trie alg\'{e}brique~I}
(Springer, Berlin, 1971; originally Publ.\ Math.\ IH\'{E}S, 1960).

\bibitem{Hartshorne1977}
R.~Hartshorne,
\textit{Algebraic Geometry}
(Springer, New York, 1977).

\bibitem{Borel1991}
A.~Borel,
\textit{Linear Algebraic Groups}, 2nd ed.\
(Springer, New York, 1991).

\bibitem{Wald1984}
R.~M.~Wald,
\textit{General Relativity}
(University of Chicago Press, Chicago, 1984).

\bibitem{Hehl1976}
F.~W.~Hehl, P.~von der Heyde, G.~D.~Kerlick and J.~M.~Nester,
Rev.\ Mod.\ Phys.\ \textbf{48}, 393 (1976).

\bibitem{Coxeter1973}
H.~S.~M.~Coxeter,
\textit{Regular Polytopes},
3rd ed.\ (Dover, New York, 1973).

\end{thebibliography}

\end{document}
